% Generated draft: requires mathematical and source review.
\documentclass[11pt]{article}
\usepackage[margin=25mm]{geometry}
\usepackage{fontspec}
\setmainfont{Latin Modern Roman}
\usepackage{amsmath,amssymb,mathtools,unicode-math}
\setmathfont{Latin Modern Math}
\usepackage{enumitem,needspace,xurl,hyperref}
\setlength{\parindent}{0pt}
\setlength{\parskip}{5pt}
\setlength{\emergencystretch}{3em}
\newcommand{\N}{\mathbb{N}}
\newcommand{\Z}{\mathbb{Z}}
\newcommand{\Q}{\mathbb{Q}}
\newcommand{\R}{\mathbb{R}}
\newcommand{\C}{\mathbb{C}}
\newcommand{\F}{\mathbb{F}}
\newcommand{\A}{\mathbb{A}}
\newcommand{\PP}{\mathbb{P}}
\newcommand{\E}{\mathbb{E}}
\newcommand{\eps}{\varepsilon}
\newcommand{\dd}{\,\mathrm{d}}
\begin{document}
\section*{UnsolvedMath ID: 30002503}
\textbf{Source problem code:} catalogue entry
\subsection{Definitions and mathematical statement}
The two usual conventions for the set of perfect squares are
\[
S_{+}=\{n^{2}:n\in\Z,\ n\geq 1\},\qquad
S_{0}=\{n^{2}:n\in\Z,\ n\geq 0\}=S_{+}\cup\{0\}.
\]
The supplied record does not specify whether zero is included.

The source question is: \emph{Is the set of perfect squares stationary for the Pisier norm?}

The record supplies neither a definition of the Pisier norm nor a definition of stationarity for that norm. In particular, it does not specify the ambient function or sequence space, the objects on which the norm acts, or the quantifiers and norm comparisons defining stationarity. Consequently, the question above is retained as a source formulation; a precise mathematical predicate cannot be supplied from this record alone. The intended definitions and the convention concerning zero require source clarification.
\subsection{Short English statement}
The question asks whether the perfect squares form a stationary set for the Pisier norm. The supplied record does not define this norm or stationarity, and does not specify whether zero is included among the squares. These omissions prevent a fully specified mathematical statement.
\subsection{Sources}
\begin{itemize}
\item Local registry: \texttt{data/unsolvedmath/problems.json}, supplied record with permanent UnsolvedMath ID \(30002503\), source problem code \texttt{catalogue entry}. The record assigns proposed year 2014 and names no proposer.
\item Supplied report citation: \emph{Dirichlet Series and Function Theory in Polydiscs} (2014), Oberwolfach Reports, DOI 10.4171/owr/2014/06, pp. 335--393. Report link and supplied literature source ``Oberwolfach Report 06/2014'': \url{https://doi.org/10.4171/owr/2014/06}. The background describes Oberwolfach Reports as documenting mathematical workshops and their open-problem discussions; no problem-specific page is supplied.
\item Supplied literature source ``Lefevre, stationary sets'': \url{https://eudml.org/doc/210549}. The record supplies no further bibliographic details or quoted definition from this reference.
\item Literature assessment dated \texttt{2026-08-21}: the supplied curation labels the problem open and reports that no published yes-or-no resolution was found. It states that later work establishes general properties and related square-frequency results without settling this question. This is supplied dated curation, not an independently verified literature assessment; the record itself requests independent verification of time-sensitive status and cited claims.
\end{itemize}
\end{document}
